\hypertarget{writing-a-garbage-collector}{%
\chapter{WRITING A GARBAGE
COLLECTOR}\label{writing-a-garbage-collector}}

C++ has RAII, but implementing a GC teaches you about the stack and
object graph.

\hypertarget{mark-and-sweep-basics}{%
\subsection{29.1 Mark-and-Sweep Basics}\label{mark-and-sweep-basics}}

\begin{enumerate}
\def\labelenumi{\arabic{enumi}.}
\tightlist
\item
  \textbf{Roots}: Pointers on the stack/globals.
\item
  \textbf{Mark}: Traverse object graph from roots, marking reachable
  objects.
\item
  \textbf{Sweep}: Iterate heap, free unmarked objects.
\end{enumerate}

\begin{lstlisting}[language={C++}]
struct GCObject {
    bool marked = false;
    virtual ~GCObject() = default;
};

class VM {
    std::vector<GCObject*> heap;
    std::vector<GCObject*> roots; // Pointers currently on stack
    
public:
    void mark() {
        for (auto* obj : roots) mark_object(obj);
    }
    
    void mark_object(GCObject* obj) {
        if (!obj || obj->marked) return;
        obj->marked = true;
        // ... traverse children ...
    }
    
    void sweep() {
        auto it = std::remove_if(heap.begin(), heap.end(), [](GCObject* obj) {
            if (!obj->marked) {
                delete obj;
                return true;
            }
            obj->marked = false; // Reset for next cycle
            return false;
        });
        heap.erase(it, heap.end());
    }
};
\end{lstlisting}

\begin{center}\rule{0.5\linewidth}{0.5pt}\end{center}
